\lysh{20722_SOLUTION}{1}
\textbf{ΛΥΣΗ}

\textbf{A)} Το άθροισμα των αποστάσεων του σημείου $K$ από τα σημεία $E(3,0)$ και $E'(-3,0)$ είναι ίσο με $10$. Επίσης $(EE')=6<10$. Συνεπώς το $K$ κινείται στην έλλειψη $C$ με εστίες τα σημεία $E$ και $E'$ και σταθερό άθροισμα $2\alpha = 10$.Είναι $2\alpha = 10 \Leftrightarrow \alpha = 5$ και $\gamma = 3$ οπότε $\beta^2 = \alpha^2 - \gamma^2 = 25-9=16$ και άρα $\beta = 4$.Η εξίσωση της $C$ είναι η $\frac{x^2}{25} + \frac{y^2}{16} = 1$.

\textbf{B)} Αρκεί να δείξουμε ότι το σύστημα
\[ \begin{cases} \frac{x^2}{25} + \frac{y^2}{16} = 1 \\ 3x+5y=25 \end{cases} \]
έχει μοναδική λύση. Από τη 2η εξίσωση έχουμε $y = \frac{25-3x}{5}$ και με αντικατάσταση στην 1η εξίσωση έχουμε
\[ \frac{x^2}{25} + \frac{\left(\frac{25-3x}{5}\right)^2}{16} = 1 \]
\[ \frac{x^2}{25} + \frac{(25-3x)^2}{25 \cdot 16} = 1 \]
\[ 16x^2 + (25-3x)^2 = 25 \cdot 16 \]
\[ 16x^2 + 625 - 150x + 9x^2 = 400 \]
\[ 25x^2 - 150x + 625 - 400 = 0 \]
\[ 25x^2 - 150x + 225 = 0 \]
\[ x^2 - 6x + 9 = 0 \]
\[ (x-3)^2 = 0 \Leftrightarrow x=3 \]
Για $x=3$ είναι $y = \frac{25-3 \cdot 3}{5} = \frac{16}{5}$, που σημαίνει ότι $C$ και $(\varepsilon)$ έχουν ένα μόνο κοινό σημείο $M(3, \frac{16}{5})$.

\textbf{Γ)} Το ότι η ευθεία $\varepsilon$ και η έλλειψη $C$ έχουν ένα μόνο κοινό σημείο, το $M$, γραφικά σημαίνει ότι η ευθεία $\varepsilon$ εφάπτεται της έλλειψης $C$ στο σημείο $M$.
Η έλλειψη $C$ έχει κορυφές τα σημεία $A(5,0)$, $A'(-5,0)$, $B(0,4)$, $B'(0,-4)$ και φαίνεται στο παρακάτω σχήμα, όπως και η εφαπτομένη $(\varepsilon)$, που εκτός από το $M$ διέρχεται και από το $(0,5)$.

% TODO: TikZ σχήμα

\textbf{Δ)} Από την ανακλαστική ιδιότητα της έλλειψης γνωρίζουμε ότι η διχοτόμος ($\delta$) της γωνίας $EME'$, είναι η κάθετη της εφαπτομένης στο σημείο $M$, όπως φαίνεται και στο παραπάνω σχήμα. Συνεπώς αναζητούμε την κάθετη στην $(\varepsilon)$ που διέρχεται από το σημείο $M$.
Η ευθεία $(\varepsilon)$ έχει συντελεστή διεύθυνσης $\lambda_\varepsilon = -\frac{3}{5}$, οπότε αφού $(\delta) \perp (\varepsilon)$, είναι $\lambda_\varepsilon \cdot \lambda_\delta = -1 \Leftrightarrow -\frac{3}{5} \cdot \lambda_\delta = -1 \Leftrightarrow \lambda_\delta = \frac{5}{3}$.
Τελικά η ζητούμενη διχοτόμος $(\delta)$ έχει εξίσωση $(\delta): y-y_M = \lambda_\delta \cdot (x-x_M) \Leftrightarrow y - \frac{16}{5} = \frac{5}{3} \cdot (x-3)$.
\end{lysh}